% Semantic definitions retained from research/machine-native-reasoning-wider/book.tex
\newcommand{\CoreTV}{d_{\mathrm{TV}}}
\newcommand{\Coreim}{\operatorname{im}}
\newcommand{\CoreSpan}{\operatorname{span}}
\newcommand{\CorePow}{\mathcal P}
\newcommand{\Coresource}[1]{\par{\small\textit{Source relationship.} #1}\par}
\newcommand{\Coreboundary}[1]{\par{\small\textit{Boundary.} #1}\par}
\newcommand{\Corequestion}[1]{\par\medskip{\large\bfseries #1}\par\smallskip}
\newtheorem{Coretheorem}{Theorem}[chapter]
\newtheorem{Corelemma}[Coretheorem]{Lemma}
\newtheorem{Coreproposition}[Coretheorem]{Proposition}
\newtheorem{Corecorollary}[Coretheorem]{Corollary}
\newtheorem{Coredefinition}[Coretheorem]{Definition}
\newtheorem{Corecounterexample}[Coretheorem]{Counterexample}

% Semantic definitions retained from monographs/komorebi-space/source/Komorebi_Space_Monograph.tex
\newcommand{\KSR}{\mathbb R}
\newcommand{\KSZ}{\mathbb Z}
\newcommand{\KSQ}{\mathbb Q}
\newcommand{\KSN}{\mathbb N}
\newcommand{\KSK}{\mathbb K}
\newcommand{\KSV}{\mathcal V}
\newcommand{\KSW}{\mathcal W}
\newcommand{\KSid}{\operatorname{id}}
\newcommand{\KSim}{\operatorname{im}}
\newcommand{\KSran}{\operatorname{ran}}
\newcommand{\KSHom}{\operatorname{Hom}}
\newcommand{\KSOb}{\operatorname{Ob}}
\newcommand{\KSExt}{\operatorname{Ext}}
\newcommand{\KSMatch}{\operatorname{Match}}
\newcommand{\KSExact}{\operatorname{Exact}}
\newcommand{\KSdist}{\operatorname{dist}}
\newcommand{\KSAut}{\operatorname{Aut}}
\newcommand{\KSStab}{\operatorname{Stab}}
\newcommand{\KSSet}{\mathbf{Set}}
\newcommand{\KSCat}{\mathbf{Cat}}
\newcommand{\KSnorm}[1]{\left\lVert#1\right\rVert}
\newcommand{\KSabs}[1]{\left\lvert#1\right\rvert}
\newcommand{\KSsource}[1]{\phantomsection\label{KS:src:#1}}
\newcommand{\KScode}[1]{\texttt{\detokenize{#1}}}
\newcommand{\KSfixedplate}[3]{\begin{figure}[H]\centering\resizebox{.96\linewidth}{!}{\input{sources/monographs/komorebi-space/source/figures/#1}}\caption{#2}\label{KS:#3}\end{figure}}
\newcommand{\KSplate}[3]{\begin{figure}[!htbp]\centering\resizebox{.96\linewidth}{!}{\input{sources/monographs/komorebi-space/source/figures/#1}}\caption{#2}\label{KS:#3}\end{figure}}
\newcommand{\KSEditionProductionModules}{162}
\newcommand{\KSEditionRoots}{4}
\newcommand{\KSEditionTargets}{166}
\newcommand{\KSEditionTheorems}{1,859}
\newcommand{\KSEditionDeclarations}{5,339}
\newcommand{\KSEditionControls}{56}
\newcommand{\KSEditionPositiveControls}{18}
\newcommand{\KSEditionNegativeControls}{35}
\newcommand{\KSEditionPolicyControls}{3}
\newcommand{\KSEditionVerifiedGroups}{16}
\newcommand{\KSEditionWrittenGroups}{0}
\newcommand{\KSEditionPrincipalGroups}{16}
\newcommand{\KSEditionSelectedMapEntries}{114}
\newcommand{\KSEditionEntrypoints}{\KSEditionRoots}
\newcommand{\KSEditionExplicitTheorems}{\KSEditionTheorems}
\newcommand{\KSEditionExportedDeclarations}{\KSEditionDeclarations}
\newcommand{\KSwritten}[1]{}
\newtheorem{KStheorem}{Theorem}[section]
\newtheorem{KSproposition}[KStheorem]{Proposition}
\newtheorem{KSlemma}[KStheorem]{Lemma}
\newtheorem{KScorollary}[KStheorem]{Corollary}
\newtheorem{KSdefinition}[KStheorem]{Definition}
\newtheorem{KSexample}[KStheorem]{Example}
\newtheorem{KSremark}[KStheorem]{Remark}

% Semantic definitions retained from monographs/grounded-abstraction/Grounded_Abstraction_v1.0.0.tex
\newcommand{\GAX}{\mathcal X}
\newcommand{\GAY}{\mathcal Y}
\newcommand{\GAZ}{\mathcal Z}
\newcommand{\GAA}{\mathcal A}
\newcommand{\GAC}{\mathcal C}
\newcommand{\GAE}{\mathcal E}
\newcommand{\GAHh}{\mathcal H}
\newcommand{\GAPp}{\mathcal P}
\newcommand{\GATV}{d_{\mathrm{TV}}}
\newcommand{\GAim}{\operatorname{im}}
\newcommand{\GAsupp}{\operatorname{supp}}
\newcommand{\GASpan}{\operatorname{span}}
\newcommand{\GAPow}{\mathcal P}
\newtheorem{GAtheorem}{Theorem}[section]
\newtheorem{GAproposition}[GAtheorem]{Proposition}
\newtheorem{GAcorollary}[GAtheorem]{Corollary}
\newtheorem{GAdefinition}[GAtheorem]{Definition}
\newtheorem{GAcounterexample}[GAtheorem]{Counterexample}

% Semantic definitions retained from research/machine-understanding/Machine_Understanding_v1.0.0.tex
\newcommand{\MUX}{\mathcal X}
\newcommand{\MUZ}{\mathcal Z}
\newcommand{\MUU}{\mathcal U}
\newcommand{\MUQ}{\mathcal Q}
\newcommand{\MUY}{\mathcal Y}
\newcommand{\MUC}{\mathcal C}
\newcommand{\MUM}{\mathcal M}
\newcommand{\MUE}{\mathbb E}
\newcommand{\MUFfive}{\mathbb F_5}
\newcommand{\MUTV}{d_{\mathrm{TV}}}
\newcommand{\MUimg}[1]{\operatorname{im}(#1)}
\newcommand{\MUEq}{\equiv_{\MUC}}
\newcommand{\MUsupp}{\operatorname{supp}}
\newcommand{\MUcode}[1]{\texttt{\small #1}}
\newtheorem{MUtheorem}{Theorem}[section]
\newtheorem{MUproposition}[MUtheorem]{Proposition}
\newtheorem{MUdefinition}[MUtheorem]{Definition}
\newtheorem{MUcounterexample}[MUtheorem]{Counterexample}
\newtheorem{MUexample}[MUtheorem]{Example}

% Semantic definitions retained from monographs/self-improving-intelligence/Self_Improving_Intelligence_v1.0.0.tex
\newcommand{\SIIX}{\mathcal X}
\newcommand{\SIIZ}{\mathcal Z}
\newcommand{\SIIU}{\mathcal U}
\newcommand{\SIIQ}{\mathcal Q}
\newcommand{\SIIY}{\mathcal Y}
\newcommand{\SIIC}{\mathcal C}
\newcommand{\SIIM}{\mathcal M}
\newcommand{\SIIE}{\mathbb E}
\newcommand{\SIIFfive}{\mathbb F_5}
\newcommand{\SIITV}{d_{\mathrm{TV}}}
\newcommand{\SIIimg}[1]{\operatorname{im}(#1)}
\newcommand{\SIIEq}{\equiv_{\SIIC}}
\newcommand{\SIIsupp}{\operatorname{supp}}
\newcommand{\SIIcode}[1]{\texttt{\small #1}}
\newcommand{\SIIB}{\mathcal B}
\newcommand{\SIIRcal}{\mathcal R}
\newcommand{\SIIGcal}{\mathcal G}
\newcommand{\SIISafe}{\mathsf{Safe}}
\newcommand{\SIIPost}{\operatorname{Post}}
\newcommand{\SIIPareto}{\operatorname{Pareto}}
\newtheorem{SIItheorem}{Theorem}[section]
\newtheorem{SIIproposition}[SIItheorem]{Proposition}
\newtheorem{SIIdefinition}[SIItheorem]{Definition}
\newtheorem{SIIcounterexample}[SIItheorem]{Counterexample}
\newtheorem{SIIexample}[SIItheorem]{Example}

% Semantic definitions retained from monographs/machine-societies/publication/Machine_Societies_v1.0.0.tex
\newcommand{\MSTV}{d_{\mathrm{TV}}}
\newcommand{\MSLaw}{\mathcal L}
\newcommand{\MSX}{\mathcal X}
\newcommand{\MSZ}{\mathcal Z}
\newcommand{\MSU}{\mathcal U}
\newcommand{\MSY}{\mathcal Y}
\newcommand{\MSHist}{\mathcal H}
\newcommand{\MSE}{\mathbb E}
\newcommand{\MSProb}{\mathbb P}
\newcommand{\MSFtwo}{\mathbb F_2}
\newcommand{\MSpush}{\pi_{\#}}
\newcommand{\MSnorm}[1]{\left\lVert#1\right\rVert}
\newcommand{\MSabs}[1]{\left\lvert#1\right\rvert}
\newcommand{\MScode}[1]{\texttt{\small #1}}
\newcommand{\MSStatus}[1]{\par\noindent\textbf{Status:} #1\par}
\newtheorem{MSproposition}{Proposition}[section]
\newtheorem{MStheorem}[MSproposition]{Theorem}
\newtheorem{MScorollary}[MSproposition]{Corollary}
\newtheorem{MSdefinition}[MSproposition]{Definition}
\newtheorem{MSexample}[MSproposition]{Example}
\newtheorem{MScounterexample}[MSproposition]{Counterexample}
\newtheorem{MShypothesis}[MSproposition]{Research hypothesis}
\newtheorem*{MSinterpretation}{Interpretation}

% Semantic definitions retained from monographs/large-signal-model/Reference_Laboratory/theory/FOUNDATIONS.tex
\newcommand{\LSMtightlist}{\setlength{\itemsep}{0pt}\setlength{\parskip}{0pt}}
\newcommand{\LSMpassthrough}[1]{#1}

% Semantic definitions retained from research/hypergraph-bayesian-refinement/Hypergraph_Lumpability_Bayesian_Refinement.tex
\newcommand{\PRIQR}{\mathbb R}
\newcommand{\PRIQC}{\mathbb C}
\newcommand{\PRIQN}{\mathbb N}
\newcommand{\PRIQcA}{\mathcal A}
\newcommand{\PRIQcB}{\mathcal B}
\newcommand{\PRIQcC}{\mathcal C}
\newcommand{\PRIQcD}{\mathcal D}
\newcommand{\PRIQcE}{\mathcal E}
\newcommand{\PRIQcF}{\mathcal F}
\newcommand{\PRIQcG}{\mathcal G}
\newcommand{\PRIQcH}{\mathcal H}
\newcommand{\PRIQcL}{\mathcal L}
\newcommand{\PRIQcP}{\mathcal P}
\newcommand{\PRIQcR}{\mathcal R}
\newcommand{\PRIQcS}{\mathcal S}
\newcommand{\PRIQcT}{\mathcal T}
\newcommand{\PRIQcV}{\mathcal V}
\newcommand{\PRIQcY}{\mathcal Y}
\newcommand{\PRIQid}{\operatorname{id}}
\newcommand{\PRIQRan}{\operatorname{Ran}}
\newcommand{\PRIQKer}{\operatorname{Ker}}
\newcommand{\PRIQSpan}{\operatorname{span}}
\newcommand{\PRIQdiag}{\operatorname{diag}}
\newcommand{\PRIQosc}{\operatorname{osc}}
\newcommand{\PRIQsupp}{\operatorname{supp}}
\newcommand{\PRIQTV}{\operatorname{TV}}
\newcommand{\PRIQLip}{\operatorname{Lip}}
\newcommand{\PRIQtr}{\operatorname{tr}}
\newcommand{\PRIQrank}{\operatorname{rank}}
\newcommand{\PRIQdd}{\,\mathrm d}
\newcommand{\PRIQnorm}[1]{\left\lVert #1\right\rVert}
\newcommand{\PRIQabs}[1]{\left\lvert #1\right\rvert}
\newcommand{\PRIQip}[2]{\left\langle #1,#2\right\rangle}
\newcommand{\PRIQfine}{\mathrm f}
\newcommand{\PRIQcoarse}{\mathrm c}
\newcommand{\PRIQBup}{\mathsf B}
\newcommand{\PRIQNrm}{\mathsf N}
\newcommand{\PRIQDef}{\mathfrak D}
\newcommand{\PRIQErr}{\mathfrak E}
\newcommand{\PRIQone}{\mathbf 1}
\newtheorem{PRIQtheorem}{Theorem}[section]
\newtheorem{PRIQproposition}[PRIQtheorem]{Proposition}
\newtheorem{PRIQlemma}[PRIQtheorem]{Lemma}
\newtheorem{PRIQcorollary}[PRIQtheorem]{Corollary}
\newtheorem{PRIQconjecture}[PRIQtheorem]{Conjecture}
\newtheorem{PRIQdefinition}[PRIQtheorem]{Definition}
\newtheorem{PRIQexample}[PRIQtheorem]{Example}
\newtheorem{PRIQalgorithm}[PRIQtheorem]{Algorithm}
\newtheorem{PRIQcertificate}[PRIQtheorem]{Numerical Certificate}
\newtheorem{PRIQboundary}[PRIQtheorem]{Claim Boundary}
\newtheorem{PRIQremark}[PRIQtheorem]{Remark}

% Semantic definitions retained from monographs/proof-path-theory/docs/Proof_Path_Theory.tex
\newcommand{\PPTcalloutbox}[2]{\par\medskip
 {\sffamily\small\bfseries\color{KomorebiTeal}#1\par}
 \smallskip #2\par\medskip}
\newcommand{\PPTPcal}{\mathcal P}
\newcommand{\PPTJcal}{\mathcal J}
\newcommand{\PPTQcal}{\mathcal Q}
\newcommand{\PPTRcal}{\mathcal R}
\newcommand{\PPTEcal}{\mathcal E}
\newcommand{\PPTVcal}{\mathcal V}
\newcommand{\PPTCcal}{\mathcal C}
\newcommand{\PPTBcal}{\mathcal B}
\newcommand{\PPTTcal}{\mathcal T}
\newcommand{\PPTFcal}{\mathcal F}
\newcommand{\PPTAnc}{\operatorname{Anc}}
\newcommand{\PPTAncL}{\operatorname{Anc}_L}
\newcommand{\PPTDesc}{\operatorname{Desc}}
\newcommand{\PPTReach}{\operatorname{Reach}}
\newcommand{\PPTCost}{\operatorname{Cost}}
\newcommand{\PPTLoss}{\operatorname{Loss}}
\newcommand{\PPTDebt}{\operatorname{Debt}}
\newcommand{\PPTargmin}{\operatorname*{arg\,min}}
\newcommand{\PPTargmax}{\operatorname*{arg\,max}}
\newcommand{\PPTopeningfigure}[3]{%
\begin{figure}[H]\centering
\includegraphics[width=\linewidth]{sources/monographs/proof-path-theory/docs/figures/#1}
\caption{#2}\label{PPT:#3}
\end{figure}}
\newcommand{\PPTstudyhead}[2]{%
 \FloatBarrier\clearpage
 {\sffamily\small\bfseries\color{KomorebiTeal}VISUAL STUDY\quad #1\par}
 \vspace{7pt}{\sffamily\fontsize{23}{27}\selectfont\bfseries\color{KomorebiInk}#2\par}\vspace{8pt}}
\newtheorem{PPTtheorem}{Theorem}[section]
\newtheorem{PPTproposition}[PPTtheorem]{Proposition}
\newtheorem{PPTlemma}[PPTtheorem]{Lemma}
\newtheorem{PPTcorollary}[PPTtheorem]{Corollary}
\newtheorem{PPTdefinition}[PPTtheorem]{Definition}
\newtheorem{PPTexample}[PPTtheorem]{Example}
\newtheorem{PPTproblem}[PPTtheorem]{Problem}
\newtheorem{PPTremark}[PPTtheorem]{Remark}
\newtheorem{PPTwarning}[PPTtheorem]{Warning}

% Semantic definitions retained from packages/constructive-repair/monograph/Constructive_Repair.tex
\newcommand{\CRkk}{\Bbbk}
\newcommand{\CRQ}{\mathbb{Q}}
\newcommand{\CRR}{\mathbb{R}}
\newcommand{\CRZ}{\mathbb{Z}}
\newcommand{\CRN}{\mathbb{N}}
\newcommand{\CRF}{\mathbb{F}}
\DeclareMathOperator{\CRim}{im}
\DeclareMathOperator{\CRcoker}{coker}
\DeclareMathOperator{\CRSpan}{span}
\DeclareMathOperator{\CRrank}{rank}
\DeclareMathOperator{\CRcost}{cost}
\DeclareMathOperator{\CRid}{id}
\DeclareMathOperator{\CRargmin}{argmin}
\newcommand{\CROb}{\mathcal{O}}
\newcommand{\CRcode}[1]{\texttt{\detokenize{#1}}}
\newcommand{\CRdecl}[1]{{\small\nolinkurl{#1}}}
\newcommand{\CRformal}[1]{\par\nobreak\begingroup\footnotesize\color{Slate}\sffamily Formal counterpart: \normalfont\nolinkurl{#1}.\par\endgroup}
\newcommand{\CRsourcescope}[1]{\par\begingroup\small\color{Slate}#1\par\endgroup}
\newcommand{\CRrecord}[3]{\noindent\begin{minipage}{\linewidth}\small\textbf{\nolinkurl{#1}}\hfill\textsf{line #2}\par\smallskip #3\end{minipage}\par\vspace{7pt}}
\newcommand{\CRmoduleheading}[2]{\section{#1}\label{CR:#2}\noindent\small Source: \nolinkurl{lean/KomorebiRepair/#1.lean}.\normalsize\par}
\newtheorem{CRtheorem}{Theorem}[section]
\newtheorem{CRproposition}[CRtheorem]{Proposition}
\newtheorem{CRlemma}[CRtheorem]{Lemma}
\newtheorem{CRcorollary}[CRtheorem]{Corollary}
\newtheorem{CRdefinition}[CRtheorem]{Definition}
\newtheorem{CRexample}[CRtheorem]{Example}
\newtheorem{CRremark}[CRtheorem]{Remark}

% Semantic definitions retained from research/stratified-relational-workspace/monograph/Stratified_Relational_Workspace_Theory.tex
\newcommand{\SRWTWorkspace}{\mathbf W}
\newcommand{\SRWTSRW}{\operatorname{SRW}}
\newcommand{\SRWTCE}{\operatorname{CE}}
\newcommand{\SRWTRisk}{\mathcal R}
\newcommand{\SRWTHol}{\operatorname{Hol}}
\newcommand{\SRWTPath}{\operatorname{Path}}
\newcommand{\SRWTRet}{\operatorname{Ret}}
\newcommand{\SRWTDep}{\operatorname{Dep}}
\newcommand{\SRWTOut}{\operatorname{Out}}
\newcommand{\SRWTCOut}{\operatorname{COut}}
\newcommand{\SRWTCIn}{\operatorname{CIn}}
\newcommand{\SRWTInstability}{\operatorname{Instability}}
\newcommand{\SRWTRec}{\operatorname{Rec}}
\newcommand{\SRWTPass}{\operatorname{Pass}}
\newcommand{\SRWTindep}{\perp\!\!\!\perp}
\newcommand{\SRWTdd}{\,\mathrm d}
\DeclareMathOperator{\SRWTran}{ran}
\DeclareMathOperator{\SRWTspec}{spec}
\DeclareMathOperator{\SRWTEnd}{End}
\DeclareMathOperator{\SRWTFix}{Fix}
\DeclareMathOperator{\SRWTComm}{Comm}
\DeclareMathOperator{\SRWTev}{ev}
\DeclareMathOperator{\SRWTim}{im}
\DeclareMathOperator{\SRWTNat}{Nat}
\DeclareMathOperator{\SRWTKL}{KL}
\DeclareMathOperator{\SRWTTV}{TV}
\DeclareMathOperator{\SRWTpr}{pr}
\DeclareMathOperator*{\SRWTargmax}{arg\,max}
\DeclareMathOperator*{\SRWTargmin}{arg\,min}
\newcommand{\SRWTcenterbox}[1]{%
\begin{center}
\fbox{\parbox{0.90\textwidth}{\centering #1}}
\end{center}
}
\newtheorem{SRWTtheorem}{Theorem}
\newtheorem{SRWTproposition}[SRWTtheorem]{Proposition}
\newtheorem{SRWTlemma}[SRWTtheorem]{Lemma}
\newtheorem{SRWTcorollary}[SRWTtheorem]{Corollary}
\newtheorem{SRWTdefinition}[SRWTtheorem]{Definition}
\newtheorem{SRWTexample}[SRWTtheorem]{Example}
\newtheorem{SRWTremark}[SRWTtheorem]{Remark}
\renewcommand{\theSRWTtheorem}{SRWT.\arabic{SRWTtheorem}}

% Semantic definitions retained from monographs/large-signal-model/Reference_Laboratory/theory/PROOF_REGISTER.tex
\newcommand{\LSMRegistertightlist}{\setlength{\itemsep}{0pt}\setlength{\parskip}{0pt}}
\newcommand{\LSMRegisterpassthrough}[1]{#1}

% Semantic definitions retained from monographs/large-signal-model/Visual_Companion/atlas.tex
\newcommand{\LSMAtlasR}{\mathbb R}
\newcommand{\LSMAtlasim}{\operatorname{im}}
\newcommand{\LSMAtlasTV}{\operatorname{TV}}
\newcommand{\LSMAtlasdiag}{\operatorname{diag}}
\newcommand{\LSMAtlasplate}[4]{\clearpage\pdfbookmark[1]{#2}{plate#1}\textcolor{teal}{\small\textsc{#1\quad #4}}\par\vspace{3pt}{\fontsize{23}{26}\selectfont #2}\par\vspace{3pt}{\color{muted}#3}\par\vspace{8pt}}
\newcommand{\LSMAtlasdiagram}[2][0.98]{\begin{center}\includegraphics[width=#1\linewidth,height=2.45in,keepaspectratio]{sources/monographs/large-signal-model/Visual_Companion/figures/#2.pdf}\end{center}\vspace{-3pt}}
\newcommand{\LSMAtlasclaim}[1]{\par\vspace{3pt}\textcolor{teal}{\textbf{The mathematical point.}} #1\par}
\newcommand{\LSMAtlasboundary}[1]{\par\vspace{3pt}{\small\textcolor{muted}{\textbf{Scope.} #1}}\par}
\newcommand{\LSMAtlasnav}[1]{\par{\small\textcolor{muted}{\textbf{Formal route.} #1}}\par}
